\section{Introduction}
\label{sec:intro}

Serving embedding-based Deep Learning recommendation models (DLRM~\cite{mudigere2022software, covington2016deep}) at scale is challenging since it is impossible to rank all items during inference. A multi-stage design is widely adopted. First, the retrieval stage narrows the item candidates to thousand scale by formulating the task as an Approximate Nearest Neighbor (ANN) search problem in vector space, identifying relevant items based on vector similarities. This is commonly built using libraries like Faiss~\cite{douze2024faiss, johnson2019billion}, RAFT~\cite{rapidsai} or a dedicated vector database system like Milvus~\cite{wang2021milvus}. Meanwhile, the retrieval stage applies feature filtering to match user attributes in multiple aspects. These filters are implemented using inverted-index mechanism. Retrieval therefore relies on the indexing and filtering services during online serving. Finally, the retrieved items are passed to downstream ranking models to generate recommendation results.

Despite wide adoption, the way of indexing and filtering in retrieval has drawbacks. First, authoring divergence of different serving components in retrieval slows down the end-to-end development, experiments and deployment. Second, isolated optimization in each component lacks co-design. Each system needs to implement redundant logic such as versioning, scheduling and batching, which makes the overall optimization fragmented~\cite{moritz2018ray}. Third, the client needs to compose multiple requests for different services and orchestrate intermediate results. This increases retrieval latency due to unnecessary data movement and transformation.
In this paper, we propose SilverTorch, a model-based retrieval system built on PyTorch. Instead of building standalone indexing services, SilverTorch defines all serving components such as ANN search and feature filtering as layers of the served model itself. Based on the unified stack, we co-design the ANN search with the feature filtering and propose an index algorithm. SilverTorch's in-model design also provides unified interface that simplifies orchestration. Client sends a single retrieval request to the model runtime that serves a SilverTorch model.

Another major challenge for serving large-scale recommendation models is compute scalability. As candidate pool grows, retrieval models need to build larger index. Meanwhile, model architectures are becoming more complex. 
Recent retrieval approaches adopt complex learned structures to replace the dot-product to represent similarities ~\cite{zhai2023revisiting}. There are also studies incorporating raw historical interaction data using transformers~\cite{zhai2024actions, zhou2018deep}. Unfortunately, existing solutions are struggling to satisfy these increasing computational demand. Existing systems adopt CPU-based ANN indexing and feature filtering~\cite{jayaram2019diskann, PinterestANN, covington2016deep, huang2020embedding}. To the best of our knowledge, no existing work studies feature filtering on GPUs. For ANN search, there are libraries implement ANN algorithms on GPUs~\cite{johnson2019billion, wang2021milvus}. However, they only support limited topk and difficult to customize for recommendation. CPU-based ANN search and filtering can scale out by adding more CPU servers to partition the index, but the cost increases linearly which quickly becomes inefficient.
SilverTorch fully utilizes GPUs, introducing model-based ANN search and bloom index filtering designed for retrieval. Both ANN search and feature filtering processes during inference are unified as tensor computation, consistent with other model serving components. We propose a co-designed indexing for ANN search and feature filtering on GPUs to further reduce GPU memory utilization and eliminate unnecessary computation. We demonstrate that SilverTorch's GPU solution is more cost-efficient.

Benefiting from SilverTorch's in-model design, we extend retrieval models by introducing an OverArch scoring layer to further learn user-item interactions. The initial returned items from ANN and filtering are re-ranked by the OverArch. For the OverArch, we pre-compute item embeddings and cache them into GPU memory to reduce online computation cost. Additionally, we enable multi-task retrieval with a aggregation layer (referred to as Value Model) to compute the overall score across multi-tasks. These initiatives enable retrieval model to pre-rank more items in the retrieval stage and improve the model consistency between retrieval and ranking stages. For ANN, we leverage Int8 quantization to save compute for pre-ranking more items in OverArch.
We summarize our contributions as follows:
\begin{itemize}%[topsep=0em, nosep, leftmargin=*]
    \item We propose SilverTorch, a unified model-based system for serving large-scale recommendation models. It unifies the serving development within PyTorch, re-defines the standalone ANN and feature filtering services with in-model tensor operators, simplifies client-side orchestration.
    \item We introduce a novel GPU-based bloom index algorithm for feature filtering and build a fused Int8 ANN kernel on GPUs. We further propose a co-designed index algorithm combining ANN search with feature filtering. This largely reduces memory utilization and eliminate unnecessary computation. To the best of our knowledge, bloom index is the first attempt to conduct feature filtering on GPUs and the first study applying it to recommendation systems.
    \item We extend existing retrieval models by introducing an OverArch scoring layer as learned similarities, as well as a multi-task retrieval with Value Model in SilverTorch. The item embeddings used for the OverArch are pre-computed and cached into SilverTorch model. These functionalities improve recommendation accuracy and enable more advanced retrieval model architectures.
    \item We evaluate SilverTorch on industry-scale datasets with a 10-million items and a 80-million items respectively. The results show SilverTorch improves serving throughput by up to $23.7\times$ compared to the state-of-the-art baselines. By introducing OverArch scoring layers and multi-task retrieval, SilverTorch achieves over 5.6\% recall improvement and is $13.35\times$ cost efficient.
\end{itemize}

% paper structure
%The paper is structured as follows. In Section \ref{sec:2}, we present background and related work. Section \ref{sec:3} shows the overall design of SilverTorch. We present the detailed design of SilverTorch model in Section \ref{sec:4} and discuss its extensibilities in Section \ref{sec:5}. Finally, we analysis the experimental results in Section \ref{sec:6} and discuss the future work in Section \ref{sec:7}. Section \ref{sec:8} concludes the paper.
